\subsection{হেক্সাডেসিমাল থেকে বাইনারি}
প্রতি অক্টাল অঙ্কে ৩ বিট, প্রতি হেক্সাডেসিমাল অঙ্কে ৪ বিট। প্রতিটি অঙ্কের বিটগুলো একই ক্রমে একত্র করো; ভগ্নাংশের শুরুর শূন্য বাদ দিও না।
\par\noindent\begin{minipage}{\linewidth}
\convExample{১}{পূর্ণসংখ্যা}
$ (A7)_{16} $-কে ২ ভিত্তিতে রূপান্তর করো।\par
\par\noindent\begin{minipage}{\linewidth}\centering
\phantomsection\label{VC052}
\begin{chThreeDiagram}\begin{tikzpicture}[x=1mm,y=1mm]
\node[vc digit] (s0) at (9.0,0) {$ \mathrm{A} $};
\draw[vc arrow] (s0.south)--(9.0,-9);
\node[vc digit] (s1) at (33.0,0) {$ 7 $};
\draw[vc arrow] (s1.south)--(33.0,-9);
\node[vc bit] (raw0) at (0,-12) {$ 1 $};
\node[vc bit] (raw1) at (6,-12) {$ 0 $};
\node[vc bit] (raw2) at (12,-12) {$ 1 $};
\node[vc bit] (raw3) at (18,-12) {$ 0 $};
\node[vc bit] (raw4) at (24,-12) {$ 0 $};
\node[vc bit] (raw5) at (30,-12) {$ 1 $};
\node[vc bit] (raw6) at (36,-12) {$ 1 $};
\node[vc bit] (raw7) at (42,-12) {$ 1 $};
\draw[vc brace] (-3.0,-16)--(21.0,-16);
\draw[vc brace] (21.0,-16)--(45.0,-16);
\end{tikzpicture}\end{chThreeDiagram}
\end{minipage}\par
\[\therefore\ (A7)_{16}=(10100111)_{2}.\]
\end{minipage}\par
\par\noindent\begin{minipage}{\linewidth}
\convExample{২}{শুধু ভগ্নাংশ}
$ (0.6D)_{16} $-কে ২ ভিত্তিতে রূপান্তর করো।\par
\par\noindent\begin{minipage}{\linewidth}\centering
\phantomsection\label{VC053}
\begin{chThreeDiagram}\begin{tikzpicture}[x=1mm,y=1mm]
\node[vc digit] (s0) at (9.0,0) {$ 0 $};
\draw[vc arrow] (s0.south)--(9.0,-9);
\node[vc digit] (s1) at (45.0,0) {$ 6 $};
\draw[vc arrow] (s1.south)--(45.0,-9);
\node[vc digit] (s2) at (69.0,0) {$ \mathrm{D} $};
\draw[vc arrow] (s2.south)--(69.0,-9);
\node[vc digit] at (27.0,0) {$.$};
\node[vc bit] (raw0) at (0,-12) {$ 0 $};
\node[vc bit] (raw1) at (6,-12) {$ 0 $};
\node[vc bit] (raw2) at (12,-12) {$ 0 $};
\node[vc bit] (raw3) at (18,-12) {$ 0 $};
\node[vc bit] (raw4) at (36.0,-12) {$ 0 $};
\node[vc bit] (raw5) at (42.0,-12) {$ 1 $};
\node[vc bit] (raw6) at (48.0,-12) {$ 1 $};
\node[vc bit] (raw7) at (54.0,-12) {$ 0 $};
\node[vc bit] (raw8) at (60.0,-12) {$ 1 $};
\node[vc bit] (raw9) at (66.0,-12) {$ 1 $};
\node[vc bit] (raw10) at (72.0,-12) {$ 0 $};
\node[vc bit] (raw11) at (78.0,-12) {$ 1 $};
\node[vc digit] at (27.0,-12) {$.$};
\draw[vc brace] (-3.0,-16)--(21.0,-16);
\draw[vc brace] (33.0,-16)--(57.0,-16);
\draw[vc brace] (57.0,-16)--(81.0,-16);
\end{tikzpicture}\end{chThreeDiagram}
\end{minipage}\par
\[\therefore\ (0.6D)_{16}=(0.01101101)_{2}.\]
পূর্ণাংশের শুরুর বা ভগ্নাংশের শেষের অপ্রয়োজনীয় ০ বাদ দিয়ে ফল লেখা হয়েছে; ভগ্নাংশের শুরুর ০ বাদ যায় না।\par
\end{minipage}\par
\par\noindent\begin{minipage}{\linewidth}
\convExample{৩}{মিশ্র সংখ্যা}
$ (4D.B8)_{16} $-কে ২ ভিত্তিতে রূপান্তর করো।\par
\par\noindent\begin{minipage}{\linewidth}\centering
\phantomsection\label{VC054}
\begin{chThreeDiagram}\begin{tikzpicture}[x=1mm,y=1mm]
\node[vc digit] (s0) at (9.0,0) {$ 4 $};
\draw[vc arrow] (s0.south)--(9.0,-9);
\node[vc digit] (s1) at (33.0,0) {$ \mathrm{D} $};
\draw[vc arrow] (s1.south)--(33.0,-9);
\node[vc digit] (s2) at (69.0,0) {$ \mathrm{B} $};
\draw[vc arrow] (s2.south)--(69.0,-9);
\node[vc digit] (s3) at (93.0,0) {$ 8 $};
\draw[vc arrow] (s3.south)--(93.0,-9);
\node[vc digit] at (51.0,0) {$.$};
\node[vc bit] (raw0) at (0,-12) {$ 0 $};
\node[vc bit] (raw1) at (6,-12) {$ 1 $};
\node[vc bit] (raw2) at (12,-12) {$ 0 $};
\node[vc bit] (raw3) at (18,-12) {$ 0 $};
\node[vc bit] (raw4) at (24,-12) {$ 1 $};
\node[vc bit] (raw5) at (30,-12) {$ 1 $};
\node[vc bit] (raw6) at (36,-12) {$ 0 $};
\node[vc bit] (raw7) at (42,-12) {$ 1 $};
\node[vc bit] (raw8) at (60.0,-12) {$ 1 $};
\node[vc bit] (raw9) at (66.0,-12) {$ 0 $};
\node[vc bit] (raw10) at (72.0,-12) {$ 1 $};
\node[vc bit] (raw11) at (78.0,-12) {$ 1 $};
\node[vc bit] (raw12) at (84.0,-12) {$ 1 $};
\node[vc bit] (raw13) at (90.0,-12) {$ 0 $};
\node[vc bit] (raw14) at (96.0,-12) {$ 0 $};
\node[vc bit] (raw15) at (102.0,-12) {$ 0 $};
\node[vc digit] at (51.0,-12) {$.$};
\draw[vc brace] (-3.0,-16)--(21.0,-16);
\draw[vc brace] (21.0,-16)--(45.0,-16);
\draw[vc brace] (57.0,-16)--(81.0,-16);
\draw[vc brace] (81.0,-16)--(105.0,-16);
\end{tikzpicture}\end{chThreeDiagram}
\end{minipage}\par
\[\therefore\ (4D.B8)_{16}=(1001101.10111)_{2}.\]
পূর্ণাংশের শুরুর বা ভগ্নাংশের শেষের অপ্রয়োজনীয় ০ বাদ দিয়ে ফল লেখা হয়েছে; ভগ্নাংশের শুরুর ০ বাদ যায় না।\par
\end{minipage}\par
